\documentclass[11pt,a4paper,uplatex,dvipdfmx]{jsarticle}

\usepackage[margin=22mm]{geometry}
\usepackage{longtable}
\usepackage{booktabs}
\usepackage{array}
\usepackage{xcolor}
\usepackage{hyperref}
\usepackage{listings}
\usepackage{enumitem}

\hypersetup{
  colorlinks=true,
  linkcolor=blue,
  urlcolor=blue
}

\lstdefinestyle{aipl}{
  basicstyle=\ttfamily\small,
  breaklines=true,
  frame=single,
  columns=fullflexible,
  keepspaces=true
}

\setlist[itemize]{leftmargin=2em}
\setlist[enumerate]{leftmargin=2em}
\renewcommand{\arraystretch}{1.2}

\title{C版AIPL ユーザーズガイド}
\author{児玉靖司（Yasushi Kodama）\\法政大学（Hosei University）}
\date{2026年5月12日}

\begin{document}
\maketitle
\tableofcontents
\newpage

\section{概要}

C版AIPLは、AIPLの \texttt{.abcl} ソースをCコードへ変換して実行するための
コード生成系である。中心となるツールは \texttt{abcl2c} であり、
実装は \texttt{src/abcl2c.ml} と \texttt{src/c\_translator.ml} にある。

C版は単一のランタイムではなく、用途別に以下の出力系を持つ。

\begin{itemize}
  \item POSIX pthread版Cコード
  \item SDL2 GUIランタイム付きCコード
  \item Xinu向けCコード
  \item 参考実装としてのPythonコード生成
\end{itemize}

本ガイドでは、主にC出力、SDL2 GUI付きC出力、Xinu向けC出力を扱う。

\section{対象実装}

\begin{longtable}{p{0.32\linewidth}p{0.58\linewidth}}
\toprule
項目 & ファイル \\
\midrule
変換器入口 & \texttt{src/abcl2c.ml} \\
C変換本体 & \texttt{src/c\_translator.ml} \\
GUIランタイム & \texttt{src/abcl\_gui\_runtime.c} \\
OCamlパーサ & \texttt{src/parser.mly}, \texttt{src/lexer.mll} \\
AST定義 & \texttt{src/ast.ml} \\
GUIサンプル & \texttt{abclc/*Gui.abcl} \\
Xinuサンプル & \texttt{abclc/*Xinu.abcl} \\
実行スクリプト & \texttt{run\_*\_\{gui,xinu\}.sh} \\
\bottomrule
\end{longtable}

\section{全体構成}

\begin{center}
\begin{tabular}{c}
\fbox{\texttt{.abcl source}} \\
$\downarrow$ \\
\fbox{\texttt{OCaml lexer/parser}} \\
$\downarrow$ \\
\fbox{\texttt{AST}} \\
$\downarrow$ \\
\fbox{\texttt{abcl2c / c\_translator.ml}} \\
$\downarrow$ \\
\fbox{POSIX C / SDL2 GUI C / Xinu C}
\end{tabular}
\end{center}

\section{abcl2cコマンド}

\texttt{abcl2c} はAIPLソースをCまたは派生ターゲットへ変換する。

\begin{lstlisting}[style=aipl]
usage: abcl2c <input.abcl> [-o <output>] [--max-msgs N] [--xinu | --python]
\end{lstlisting}

\begin{longtable}{p{0.30\linewidth}p{0.60\linewidth}}
\toprule
オプション & 意味 \\
\midrule
\texttt{-o <output>} & 出力ファイル名を指定する。省略時は入力名から拡張子を変えて出力する。 \\
\texttt{--max-msgs N} & 実行時に処理するメッセージ数上限。0で上限無効。 \\
\texttt{--xinu} & Xinu向けCコードを生成する。 \\
\texttt{--python} & 参考用のPythonコードを生成する。 \\
\bottomrule
\end{longtable}

\section{POSIX pthread版C}

標準のC出力は、POSIX pthreadを使うアクターランタイムを内包する。
各AIPLアクターはC側の \texttt{object\_t} として管理され、各オブジェクトに
メールボックスとpthreadが割り当てられる。

\subsection{主なランタイム構造}

\begin{itemize}
  \item \texttt{value\_t}: \texttt{int}, \texttt{float}, \texttt{string}, object idを表す値。
  \item \texttt{message\_t}: sender、receiver、method、引数列を持つメッセージ。
  \item \texttt{mailbox\_t}: mutexとcondition variableを持つリングバッファ型メールボックス。
  \item \texttt{object\_t}: class id、fields、mailbox、pthreadを持つアクター。
  \item \texttt{enqueue}: 受信側メールボックスへ非同期メッセージを投函する。
  \item \texttt{actor\_main}: 各アクターのメッセージ処理ループ。
\end{itemize}

\subsection{生成される処理}

変換後のCコードでは、AIPLクラスごとに以下が生成される。

\begin{itemize}
  \item class idの定義
  \item フィールド番号のenum
  \item フィールド初期化関数
  \item メソッド関数
  \item method名によるdispatch関数
  \item グローバルactor生成と起動処理
\end{itemize}

\subsection{通常のビルド例}

\begin{lstlisting}[style=aipl,language=bash]
cd /Users/yaskodama/local-genai-chatgpt-ga/aios-claude
dune build src/abcl2c.exe
./_build/default/src/abcl2c.exe abclc/PingPong.abcl -o /tmp/pingpong.c --max-msgs 12
cc -O2 -Wall -pthread -o /tmp/pingpong /tmp/pingpong.c -lm
/tmp/pingpong
\end{lstlisting}

\section{対応するAIPL機能}

C版が主に扱うAIPL機能は以下である。

\begin{longtable}{p{0.30\linewidth}p{0.60\linewidth}}
\toprule
機能 & 対応状況 \\
\midrule
\texttt{class} / \texttt{method} & C関数とdispatch関数へ変換される。 \\
\texttt{var} フィールド & C配列 \texttt{fields[]} の要素として保持される。 \\
\texttt{new Class(args)} & C側でオブジェクトを確保し、\texttt{init} メッセージを投函する。 \\
\texttt{send actor.method(args)} & \texttt{enqueue} に変換される。 \\
\texttt{send!} & 通常の \texttt{send} と同様にC側投函へ変換される。 \\
\texttt{self} & 現在の \texttt{self\_id}。 \\
\texttt{sender} & メッセージ送信元ID。 \\
\texttt{if} / \texttt{while} & Cの制御構文へ変換される。 \\
算術・比較 & \texttt{v\_binop} で処理される。 \\
\texttt{print} & \texttt{v\_print} で標準出力へ出す。 \\
外部組み込み & Cの \texttt{extern value\_t name(int,value\_t*)} としてリンクする。 \\
\bottomrule
\end{longtable}

\section{未対応または制限のある機能}

\begin{itemize}
  \item \texttt{become} はC出力ではコメント \texttt{/* become unsupported */} になる。
  \item \texttt{select} はC出力ではコメント \texttt{/* select unsupported */} になる。
  \item \texttt{now} / \texttt{future} / \texttt{await} はC出力の中心機能ではない。
  \item Python版AIPLの型注釈、所有権、線形型、transient cast、モデル検査機能はC変換対象ではない。
  \item 標準ランタイムの固定上限として、\texttt{MAX\_OBJECTS}, \texttt{MAX\_MAILBOX}, \texttt{MAX\_FIELDS}, \texttt{MAX\_ARGS} がある。
  \item non-objectなトップレベル変数は標準C出力では限定的な扱いになる。
  \item Xinu版はPOSIX版よりさらに小さい固定上限を持つ。
\end{itemize}

\section{SDL2 GUI付きC版}

GUI付きC版は、通常のC出力に \texttt{src/abcl\_gui\_runtime.c} をリンクして使う。
SDL2を使い、線描画、ボタン、スライダー、哲学者、bounded buffer、災害/ドローン風
可視化などを提供する。

\subsection{代表的なGUI組み込み}

\begin{longtable}{p{0.34\linewidth}p{0.56\linewidth}}
\toprule
関数 & 概要 \\
\midrule
\texttt{gui\_open(w,h,title)} & GUIウィンドウを開く準備をする。 \\
\texttt{gui\_set\_line(...)} & 線分を描画状態へ設定する。 \\
\texttt{gui\_add\_button(...)} & GUIボタンを追加する。 \\
\texttt{gui\_register\_ticker(actor)} & 定期的に動かすactorを登録する。 \\
\texttt{gui\_run()} & SDLイベントループを実行する。 \\
\texttt{gui\_dining\_init(N)} & 食事する哲学者の配置を初期化する。 \\
\texttt{gui\_set\_phil(i,state)} & 哲学者状態を更新する。 \\
\texttt{gui\_set\_fork\_held/free} & フォーク状態を更新する。 \\
\texttt{gui\_buf\_setup(...)} & bounded buffer表示を初期化する。 \\
\texttt{gui\_add\_slider(...)} & スライダーを追加する。 \\
\texttt{gui\_slider\_value(i)} & スライダー値を読む。 \\
\bottomrule
\end{longtable}

\subsection{GUIビルド例}

\begin{lstlisting}[style=aipl,language=bash]
cd /Users/yaskodama/local-genai-chatgpt-ga/aios-claude
dune build src/abcl2c.exe
./_build/default/src/abcl2c.exe abclc/Rotate4LinesGui.abcl -o /tmp/r4l.c --max-msgs 0
cc -O2 -Wall -pthread $(pkg-config --cflags sdl2) \
   -o /tmp/r4l /tmp/r4l.c src/abcl_gui_runtime.c \
   $(pkg-config --libs sdl2) -lm
/tmp/r4l
\end{lstlisting}

既存スクリプトとして、以下が用意されている。

\begin{itemize}
  \item \texttt{run\_r4l\_gui.sh}
  \item \texttt{run\_p5\_gui.sh}
  \item \texttt{run\_bb\_gui.sh}
  \item \texttt{run\_disaster\_gui.sh}
\end{itemize}

\section{Xinu向けC版}

\texttt{--xinu} を指定すると、Xinu向けCコードを生成する。POSIX pthreadの代わりに、
Xinuの \texttt{thread}, \texttt{semaphore}, \texttt{kprintf},
\texttt{create}, \texttt{ready} を使うランタイムになる。

\subsection{Xinu版の特徴}

\begin{itemize}
  \item \texttt{pthread} ではなくXinu threadを使う。
  \item \texttt{enqueue} はXinu内部名との衝突を避けて \texttt{abcl\_enqueue} として実装される。
  \item エントリポイントは \texttt{thread aipl\_main(void)} として生成される。
  \item QEMU上のXinu kernelへ組み込んで実行する想定である。
\end{itemize}

\subsection{Xinu実行例}

\begin{lstlisting}[style=aipl,language=bash]
cd /Users/yaskodama/local-genai-chatgpt-ga/aios-claude
dune build src/abcl2c.exe
./_build/default/src/abcl2c.exe abclc/PingPong.abcl \
    -o /tmp/pp_xinu.c --xinu --max-msgs 12
\end{lstlisting}

生成したCファイルをXinuの \texttt{apps/} に配置し、Xinu kernelをビルドして
QEMUで実行する。既存スクリプトとして以下がある。

\begin{itemize}
  \item \texttt{run\_pingpong\_xinu.sh}
  \item \texttt{run\_r4l\_xinu.sh}
  \item \texttt{run\_p5\_xinu.sh}
  \item \texttt{run\_bb\_xinu.sh}
\end{itemize}

\section{サンプル分類}

\begin{longtable}{p{0.34\linewidth}p{0.56\linewidth}}
\toprule
サンプル & 用途 \\
\midrule
\texttt{abclc/PingPong.abcl} & 基本的なアクター間メッセージ。 \\
\texttt{abclc/Rotate4LinesGui.abcl} & SDL2 GUIで回転線を描画。 \\
\texttt{abclc/Philosophers5Gui.abcl} & GUI付き食事する哲学者。 \\
\texttt{abclc/BoundedBufferGui.abcl} & GUI付きbounded buffer。 \\
\texttt{abclc/DisasterReturnGui.abcl} & 災害/帰還支援風GUIサンプル。 \\
\texttt{abclc/Rotate4LinesXinu.abcl} & Xinu向け回転線サンプル。 \\
\texttt{abclc/Philosophers5Xinu.abcl} & Xinu向け哲学者サンプル。 \\
\texttt{abclc/BoundedBufferXinu.abcl} & Xinu向けbounded buffer。 \\
\bottomrule
\end{longtable}

\section{テスト}

\texttt{abclc/\_smoke\_test.sh} は、OCaml REPLで動かすサンプルと、
\texttt{abcl2c} で変換すべきGUI/Python/Xinu系サンプルを分けて検査する。
GUI系は \texttt{abcl2c} でC生成後、\texttt{cc} とSDL2でコンパイルする。

\begin{lstlisting}[style=aipl,language=bash]
cd /Users/yaskodama/local-genai-chatgpt-ga/aios-claude/abclc
./_smoke_test.sh
\end{lstlisting}

\section{まとめ}

C版AIPLは、AIPLプログラムをネイティブCランタイムへ落とし込むための実験的な
コード生成系である。POSIX pthread版では軽量なアクター実行を確認でき、
SDL2 GUIランタイムをリンクすれば可視化デモを構築できる。さらにXinu向け出力により、
教育用OSや組込み風環境でAIPLアクターを動かす実験も可能である。

\end{document}
